\section*{Hoofdstuk \ref{ch:b3:ocean-dynamics} --- Oceaandynamica}

\begin{solution}{exo:b3:ocean-dynamics:1}
$f = \qty{9.37e-5}{s^{-1}}$; $\tau/(\rho f) = 0.12/(1025\times\num{9.37e-5})
= \qty{1.25}{m^2/s}$, naar rechts van de wind: naar het zuiden. Over
\qty{1000}{km}, \qty{1.25e6}{m^3/s}, \qty{1.25}{\sverdrup}.
\end{solution}

\begin{solution}{exo:b3:ocean-dynamics:2}
$f = \qty{1.031e-4}{s^{-1}}$; $D_E = \pi\sqrt{2\times0.05/\num{1.031e-4}}
= \qty{98}{m}$. De stroom daar is $\mathrm e^{-\pi} = \num{0.043}$ van de
oppervlaktestroom, en stroomt in de tegengestelde richting.
\end{solution}

\begin{solution}{exo:b3:ocean-dynamics:3}
$\beta = 2\Omega\cos\ang{30}/a = \qty{1.98e-11}{m^{-1}.s^{-1}}$; $V =
\num{-1e-7}/(1025\times\num{1.98e-11}) = \qty{-4.92}{m^2/s}$; over
\qty{5000}{km}, \qty{-24.6}{\sverdrup}: naar het zuiden.
\end{solution}

\begin{solution}{exo:b3:ocean-dynamics:4}
De twee componenten komen na $1/(1/12 - 1/12.4206) =
\qty{354}{h}$ weer in fase, \num{14.77}\text{ dagen}, een halve synodische maand.
\end{solution}

\begin{solution}{exo:b3:ocean-dynamics:5}
Op \ang{40}~NB geldt $f = \qty{9.37e-5}{s^{-1}}$, $\beta =
\qty{1.75e-11}{m^{-1}.s^{-1}}$, $\tau_x = \qty{0.05}{Pa}$ en
$\partial\tau_x/\partial y = 0.1/\num{1.11e6} = \qty{9.0e-8}{Pa/m}$. De
eerste term is $-\num{9.0e-8}/(1025\times\num{9.37e-5}) =
\qty{-9.38e-7}{m/s}$; de tweede, $0.05\times\num{1.75e-11}/(1025\times
(\num{9.37e-5})^2) = \qty{+9.7e-8}{m/s}$. $w_E = \qty{-8.4e-7}{m/s}$,
\qty{-26.5}{m/yr}: het water wordt naar beneden geduwd.
\end{solution}

\begin{solution}{exo:b3:ocean-dynamics:6}
$\delta_S = \num{2e-6}/\num{2e-11} = \qty{100}{km}$; snelheid $=
\num{30e6}/(\num{1e5}\times800) = \qty{0.375}{m/s}$.
\end{solution}

\begin{solution}{exo:b3:ocean-dynamics:7}
$\beta = \qty{1.98e-11}{m^{-1}.s^{-1}}$; $\delta_M = (10^4/\num{1.98e-11})^{1/3}
= \qty{79.6}{km}$, en de straalstroom is $3.628\,\delta_M = \qty{289}{km}$ breed.
Bij $A_H$ die acht keer kleiner is, worden beide gehalveerd, omdat $8^{1/3} = 2$.
\end{solution}

\begin{solution}{exo:b3:ocean-dynamics:8}
$f = \qty{8.57e-5}{s^{-1}}$; $v = (21.3 - 14.5)/(\num{8.57e-5}\times\num{8e4})
= \qty{0.99}{m/s}$, positief met $B$ rechts: kijkend van $A$ naar $B$,
naar het oosten, stroomt de stroom naar het noorden. Het zeeoppervlak staat
$6.8/9.81 = \qty{0.69}{m}$ hoger bij $B$; de waterkolom bij $B$ is de
lichtere, warmere, aan de rechterzijde van de stroom.
\end{solution}

\begin{solution}{exo:b3:ocean-dynamics:9}
Springtijslag $1.5K(1 + 0.46) = \qty{0.78}{m}$; doodtijslag $1.5K(1 - 0.46)
= \qty{0.29}{m}$; verhouding $1.46/0.54 = 2.7$.
\end{solution}

\begin{solution}{exo:b3:ocean-dynamics:10}
Zoals in \cref{thm:b3:ocean-dynamics:kelvin}, $u = \pm(g/c)\eta$ voor $F(x \mp
ct)$, en $fu = -g\,\partial\eta/\partial y$ wordt $\beta_0y(\pm g/c)G =
-gG'$, dus $G = \exp(\mp\beta_0y^2/2c)$. Alleen het bovenste teken, de
oostwaarts lopende golf, is aan beide zijden van de evenaar begrensd. $\beta_0 =
\qty{2.29e-11}{m^{-1}.s^{-1}}$, $\sqrt{2.5/\num{2.29e-11}} = \qty{330}{km}$.
\end{solution}

\begin{solution}{exo:b3:ocean-dynamics:11}
(a)~Integreer $\beta v = f\,\partial w/\partial z$ vanaf de bodem, $w = 0$,
tot de top van de abyssale laag, $w = W_0$: $\beta V = fW_0$, en $f/\beta =
a\tan\varphi$. (b)~$V = \num{6.371e6}\times0.577\times10^{-7} =
\qty{0.368}{m^2/s}$, \qty{1.84}{\sverdrup} over \qty{5000}{km}, naar de pool.
(c)~Van een diepe westelijke randstroom die gevoed wordt door de bronnen op
hoge breedte; het inwendige stroomt naar de bronnen toe, de randstroom er
vandaan.
\end{solution}

\begin{solution}{exo:b3:ocean-dynamics:12}
(a)~Met $z$ naar boven vanaf de bodem, $A_zW'' = \mathrm if(W - W_g)$, waarbij de
drukgradiënt die van de geostrofische stroming boven is. $W - W_g =
C\mathrm e^{-(1+\mathrm i)z/\delta}$ is naar boven begrensd, en $W(0) = 0$ geeft
$C = -W_g$. (b)~Het transporttekort is $\int_0^\infty(W - W_g)\,\mathrm dz
= -W_g\delta/(1+\mathrm i) = -\tfrac12\delta(1 - \mathrm i)W_g$, dat is
$-\tfrac12\delta(u_g + v_g,\ v_g - u_g)$. Omdat de geostrofische stroming
niet-divergent is, geldt $w = -\nabla\cdot(\text{tekort}) = \tfrac12\delta(\partial_x
v_g - \partial_y u_g) = \tfrac12\delta\zeta_g$. (c)~$\delta =
\sqrt{0.02/10^{-4}} = \qty{14.1}{m}$; $r = 10^{-4}\times14.1/8000 =
\qty{1.77e-7}{s^{-1}}$, een afremtijd van \num{65}\text{ dagen}.
\end{solution}

\begin{solution}{pb:b3:ocean-dynamics:1}
\textbf{1.} $\tau_x = \qty{-0.1}{Pa}$ (oostenwind), 0 en \qty{+0.1}{Pa}
(westenwind).

\textbf{2.} $f = \num{4.99e-5}$, \num{8.37e-5} en
\qty{1.117e-4}{s^{-1}}.

\textbf{3.} $V_E = -\tau_x/(\rho f)$: \qty{+1.96}{m^2/s} (noord) bij
\ang{20}~NB en \qty{-0.87}{m^2/s} (zuid) bij \ang{50}~NB. Ze convergeren.

\textbf{4.} $D_E = \pi\sqrt{2\times0.05/\num{8.37e-5}} = \qty{109}{m}$.

\textbf{5.} $\tau_x = 0$ bij \ang{35}~NB, dus $w_E = -(1/\rho f)\,
\partial\tau_x/\partial y = -\pi\tau_0/(\rho fL) = \qty{-1.10e-6}{m/s}$,
\qty{-34.7}{m/yr}, naar beneden.

\textbf{6.} $500/34.7 = \num{14.4}\text{ jaar}$.

\textbf{7.} $L = a\times\pi/6 = \qty{3336}{km}$; $\beta =
\qty{1.875e-11}{m^{-1}.s^{-1}}$.

\textbf{8.} $\operatorname{curl}_z\vec\tau = -\pi\tau_0/L =
\qty{-9.42e-8}{Pa/m}$.

\textbf{9.} $V = \num{-9.42e-8}/(1025\times\num{1.875e-11}) =
\qty{-4.90}{m^2/s}$, naar het zuiden.

\textbf{10.} $4.90\times\num{6e6} = \qty{2.94e7}{m^3/s}$,
\qty{29.4}{\sverdrup} naar het zuiden.

\textbf{11.} $\qty{4.9}{mm/s}$; \qty{1000}{km} in \qty{2.0e8}{s},
\num{6.5}~jaar.

\textbf{12.} $fw_E/\beta = \num{8.37e-5}\times(\num{-1.10e-6})/
\num{1.875e-11} = \qty{-4.90}{m^2/s}$: hetzelfde, omdat het
Ekmantransport verdwijnt waar $\tau_x = 0$. Een kolom die met de snelheid $w_E$
samengeperst, behoudt zijn \omterm{def:b3:atmospheric-dynamics:pv}{potentiële vorticiteit} $f/h$ alleen door naar een kleinere $f$ te schuiven, naar het zuiden,
met de snelheid $\beta V_g = fw_E$.

\textbf{13.} Een grenslaag kan alleen bestaan als de correctie op
het binnenste deel van de stroming afneemt vanaf de wand; met $\beta > 0$ is dat
$\mathrm e^{-x/\delta_S}$, die naar het oosten afneemt, alleen vanaf de westelijke
wand. In termen van vorticiteit zou wrijving in een noordwaartse stroming tegen een
oostelijke wand positieve vorticiteit wegnemen, het verkeerde teken.

\textbf{14.} $\delta_S = \num{e-6}/\num{1.875e-11} = \qty{53.3}{km}$.

\textbf{15.} $A = \qty{4.90}{m^2/s}$, $X/\delta_S = 112.5$:
$4.90\times111.5 = \qty{546}{m^2/s}$, \qty{0.55}{m/s} door
\qty{1000}{m}.

\textbf{16.} Waar $(X/\delta_S)\mathrm e^{-x/\delta_S} = 1$: $x =
\delta_S\ln(X/\delta_S) = \qty{252}{km}$.

\textbf{17.} $A(X - \delta_S - x_0) = 4.90\times(6000 - 53 - 252)\times
10^3 = \qty{27.9}{\sverdrup}$. De stroming in het binnenste deel gaat in de
strook naast de wand nog steeds naar het zuiden; het tekort, $A(x_0 + \delta_S) =
\qty{1.5}{\sverdrup}$, is het transport in het binnenste deel dat door
die strook had moeten stromen.

\textbf{18.} $\delta_M = (\num{5e3}/\num{1.875e-11})^{1/3} = \qty{64.4}{km}$;
straal $3.628\times64.4 = \qty{233}{km}$.

\textbf{19.} $1.163\times29.4 - A\times\qty{233}{km} = 34.2 - 1.1 =
\qty{33.0}{\sverdrup}$; het overschot, \qty{16}{\percent} van het transport in het
binnenste deel keert naar het zuiden terug in de tegenstroom net ten oosten van de straal.

\textbf{20.} $c = \sqrt{9.81\times4000} = \qty{198}{m/s}$; $L_R =
198/\num{8.37e-5} = \qty{2370}{km}$.

\textbf{21.} $198\times12.4206\times3600 = \qty{8860}{km}$.

\textbf{22.} Omtrek $2(6000 + 3336) = \qty{18670}{km}$, \qty{26.2}{h},
2.1 M$_2$-perioden.

\textbf{23.} De cirkel is $2\pi a\cos\ang{35} = \qty{32790}{km}$, die wordt
in één lunaire dag doorlopen, \qty{24.84}{h}: \qty{367}{m/s}, wat $c^2/g =
\qty{13.7}{km}$ water vereist.

\textbf{24.} Met een Rossbystraal die vergelijkbaar is met de breedte van het bekken
en een M$_2$-golflengte die vergelijkbaar is met de lengte ervan, ontstaat ongeveer één \omterm{def:b3:ocean-dynamics:amphidrome}{amfidromisch
punt}, nabij het midden van het bekken; het hoogwater draait tegen de klok in
eromheen en houdt de kust aan zijn rechterkant.

\textbf{25.} $\tfrac34\times0.357\times\cos^2\ang{35} = \qty{0.18}{m}$.

\textbf{26.} De westelijke randstroom voert naar het noorden de
\textbf{\qty[text-series-to-math=true]{29}{\sverdrup}} af die het binnenste deel
naar het zuiden voert (\qty{29.4}{\sverdrup} volgens de berekening).
\end{solution}
